\documentstyle[%


righttag%


,11pt%


,amssymb%


,russian%               remove in Eng.


,izvru%                 Set izven% in Eng.


]{amsart}





\newcommand{\sh}{\operatorname{sh}}


\newcommand{\cth}{\operatorname{cth}}





\newcommand{\tts}[1]{}





\theoremstyle{plain}


\theorembodyfont{\it}


\newtheorem{lemnonum}{\hskip\parindent Лемма}


\renewcommand{\thelemnonum}{}








%\hoffset=-15mm%                  For Eng.


\setcounter{page}{39}


\god{2005}


\nomer{8~(519)} %                        Remove in Eng.


%\nomer{Vol.\,49, No.\,8}%               Set in Eng.


%\str{\pageref{begin}--\pageref{end}}%  Set in Eng.


\udk{517.553}





\author{\it А.М.\,Кытманов, З.Е.\,Потапова}


% NB:   ^^ remove in Eng.


\title{Формулы для нахождения степенных сумм 


корней систем мероморфных функций}


\thanks{Работа первого автора выполнена при финансовой поддержке Российского 


фонда фундаментальных исследований, грант \No\,02-01-00167; работа второго


--- при поддержке Совета по грантам при Президенте Российской Федерации


(\No\,НШ~1212-2003.1).}





\date{}


%\pagestyle{myheadings}%         Set in Eng.


\pagestyle{plain} %              Remove in Eng.


\begin{document}


%\thispagestyle{plain}%          Set in Eng.


\maketitle


\label{begin}








\section*{Введение}





Рассмотрим систему функций $f_1(z), f_2(z),\dotsc, f_n(z)$, голоморфных в


окрестности точки $0\in\Bbb C^n$, $z=(z_1,z_2,\dotsc, z_n)$, имеющих


следующий вид:


\begin{equation}\label{f1}


f_j(z)=z^{\beta^j}+Q_j(z),\quad j=1,2,\ldots, n,


\end{equation}


где $\beta^j =(\beta_1^j,\beta_2^j,\dotsc,\beta_n^j)$ --- мультииндекс с


целыми неотрицательными координатами, $z^{\beta^j}=z_1^{\beta_1^j}\cdot


z_2^{\beta_2^j}\dotsb z_n^{\beta_n^j}$ и $\|\beta^j\|=\beta_1^j+\beta_2^j+


\dotsb +\beta_n^j=k_j$, $j=1,2,\dotsc, n$. Функции $Q_j$ разлагаются в


окрестности нуля в ряд Тейлора, сходящийся  абсолютно и равномерно,


\begin{equation}\label{f2}


Q_j(z)=\sum_{\|\alpha\|\geqslant 0}a_\alpha^j z^\alpha,


\end{equation}


где $\alpha=(\alpha_1,\alpha_2,\dotsc,\alpha_n)$,


$\alpha_j\geqslant 0$, $\alpha_j\in\Bbb Z$,


а


$z^\alpha=z_1^{\alpha_1}\cdot z_2^{\alpha_2}\dotsb z_n^{\alpha_n}$.





В дальнейшем будем считать, что степени всех мономов (по совокупности 


переменных), входящих в $Q_j$, строго больше $k_j$,


$j=1,2,\dotsc, n$ ($\|\alpha\|=


\alpha_1+\alpha_2+\dotsb +\alpha_n>k_j$).





Рассмотрим циклы $\gamma(r)=\gamma(r_1,r_2,\dotsc, r_n)$, являющиеся 


остовами поликругов,


$$


\gamma(r)=\{z\in\Bbb C^n:|z_s|=r_s,\ s=1,2,\dotsc,n\},


\quad r_1>0,\dotsc,r_n>0.


$$





При достаточно малых $r_j$ циклы $\gamma(r)$ лежат в области голоморфности


функций $f_j$, поэтому ряды


$$


\sum_{\|\alpha\|\geqslant 0}|a^j_\alpha |r_1^{\alpha_1}


\dotsb r^{\alpha_n}_n, \quad j=1,2,\dotsc, n,


$$


сходятся. Тогда на цикле 


$\gamma(tr)=\gamma(tr_1,tr_2,\dotsc,tr_n)$, $t>0$,


имеем


$$


|z|^{\beta^j}=t^{k_j}\cdot r_1^{\beta_1^j}\cdot r_2^{\beta_2^j}\dotsb


r_n^{\beta_n^j}=t^{k_j}\cdot r^{\beta^j},


$$


а


$$


|Q_j(z)|=\bigg|\sum_{\|\alpha\|\geqslant 0}a^j_\alpha z^\alpha\bigg|


\leqslant\sum_{\|\alpha\|\geqslant 0}t^{\|\alpha\|}


|a^j_\alpha|r^\alpha\leqslant


t^{k_j+1}\sum_{\|\alpha\|\geqslant 0} |a_\alpha^j|r^\alpha.


$$


Поэтому при достаточно малых положительных $t$ на цикле $\gamma(tr)$


выполняются неравенства


\begin{equation}\label{f3}


|z|^{\beta^j}>|Q_j(z)|,\quad j=1,2,\dotsc, n,


\end{equation}


таким образом,


$$


f_j(z)\neq 0\ \ \text{на}\ \ \gamma(tr),\quad j=1,2,\dotsc,n.


$$


В частности, в некоторой достаточно малой окрестности нуля система


\begin{equation}\label{f4}


f_j(z)=0,\quad j=1,\dotsc,n,


\end{equation}


может иметь корни только на координатных плоскостях $\{z: z_s=0\}$, 


$s=1,2,\dotsc,n$ (по принципу максимума для голоморфных функций).





Из (\ref{f3}) следует, что при достаточно малых $r_j$ определены интегралы


$$


\int\limits_{\gamma(r)}\frac 1{z^{\beta+I}}\frac{df}{f}=


\int\limits_{\gamma(r_1,r_2,\dotsc,r_n)}


\frac{1}{z_1^{\beta_1+1}\cdot z_2^{\beta_2+1}\dotsb


z_n^{\beta_n+1}}\frac{df_1}{f_1}\land\frac{df_2}{f_2}\land\dotsb\land


\frac{df_n}{f_n},


$$


где $\beta_1\geqslant 0,\, \beta_2\geqslant 0,\dotsc, \beta_n\geqslant 0$,


$\beta_j\in\Bbb Z$, $I=(1,1,\dotsc, 1)$. По теореме Коши--Пуанкаре эти


интегралы не зависят от $(r_1,\dotsc,r_n)$. Обозначим


$$


J_\beta=\frac 1{(2\pi i)^n}\int_{\gamma(r)}\frac


1{z^{\beta+I}}\frac{df}{f}.


$$





Наша цель состоит в решении следующих задач: 





1) в вычислении этого интеграла через коэффициенты функций


$Q_j(z)$ (теорема 1);





2) при дополнительных условиях на функции $f_j$


(если $f_j$ ---  целые или мероморфные функции) в выяснении его


связи со степенными суммами корней (и полюсов) системы


(\ref{f4}) (теоремы 2 и 3).





Если $f_j$ --- полиномы (т.\,е. система (\ref{f4}) есть система 


алгебраических уравнений), то известны формулы для вычисления 


степенных сумм корней этой системы через коэффициенты $f_j$ 


(напр., \cite{1}, \cite{2}). На этих формулах


основан модифицированный метод исключения неизвестных, предложенный


Л.А.\,Айзенбергом в \cite{1} и развитый затем в \cite{3}. Его компьютерная


реализация дана в \cite{4}, \cite{5}.





В этих работах рассматривались степенные суммы корней в положительной 


степени. Если в качестве $f_j$ взять целые (или мероморфные)


функции, то система (\ref{f4}) может иметь бесконечное число корней (или


полюсов) и, следовательно, степенные суммы в положительной степени могут


быть не определены. В данной работе для степенных сумм


корней в отрицательной степени получены конечные формулы для их


нахождения.








\section{Основные результаты}





Пусть для системы (\ref{f4}) выполнены условия, сформулированные


во введении.








\begin{thm}\label{t1}


При сделанных предположениях для функции $f_j$ вида $(\ref{f1})$,


$(\ref{f2})$


\begin{gather*}


J_\beta=\sum_{\|\alpha\|\leqslant\|\beta\|+\min(n,k_1+\dotsb+k_n)}


\frac{(-1)^{\|\alpha\|}}{(\beta+(\alpha_1+1)\beta^1+\dotsb


+(\alpha_n+1)\beta^n)!}\frac{\partial^k(\Delta\cdot Q^\alpha)}


{\partial z^{\beta+(\alpha_1+1)\beta^1+\dotsc+(\alpha_n+1)\beta^n}}


\bigg|_{z=0}= \\


%


=\sum_{\|\alpha\|\leqslant\|\beta\|+\min(n,k_1+\dotsb+k_n)}(-1)^{\|\alpha\|}


\frak{M}\bigg[\frac{\Delta\cdot Q^\alpha}{z^{\beta+(\alpha_1+1)\beta^1+


\dotsb+(\alpha_n+1)\beta^n}}\bigg],


\end{gather*}


где $k=\|\beta+(\alpha_1+1)\beta^1+\dotsb+(\alpha_n+1)\beta^n\|$,


$\beta!=\beta_1!\cdot\beta_2!\dotsb\beta_n!$, $\Delta$ --- якобиан системы


$(\ref{f4})$,


$Q^\alpha=Q_1^{\alpha_1}\cdot Q_2^{\alpha_2}\dotsb Q_n^{\alpha_n}$,


$\frac{\partial^{\|\beta\|}}{\partial z^\beta}=\frac{\partial^{\|\beta\|}}


{\partial z_1^{\beta_1}\partial z_2^{\beta_2}\dotsb\partial z_n^{\beta_n}}$


и, наконец, $\frak M$ --- линейный функционал, сопоставляющий ряду Лорана


его свободный член.


\end{thm}








{\bf Доказательство.}


Поскольку для $\gamma(r)$ выполнены неравенства (\ref{f3}), то


$$


\frac 1{f_j}=\frac 1{z^{\beta^j}+Q_j(z)}=


\sum_{s=0}^\infty (-1)^s\frac{(Q_j(z))^s}{z^{(s+1)\beta^j}}


$$


и ряд сходится абсолютно и равномерно на $\gamma(r)$,


$j=1,2,\dotsc, n$. Поэтому


$$


\frac 1{z^{\beta+I}}\frac{df}{f}=\frac{\Delta\, dz}{z^{\beta+I}(z^{\beta_1}


+Q_1)(z^{\beta_2}+Q_2)\dotsb (z^{\beta_n}+Q_n)}=


\frac{\Delta\, dz}{z^{\beta+I}}\sum_{\|\alpha\|\geqslant 0}(-1)^{\|\alpha\|}


\frac{Q_1^{\alpha_1}\cdot Q_2^{\alpha_2}\dotsb


Q_n^{\alpha_n}}{z^{(\alpha_1+1)\beta^1


+\dotsb+(\alpha_n+1)\beta^n}},


$$


где $\Delta=\det \big\|\frac{\partial f_j}{\partial z_k}\big\|_{j,k=1}^n$


--- якобиан системы (\ref{f4}), а $dz=dz_1\land dz_2\land\dotsb\land dz_n$.





Следовательно, на $\gamma(r)$ получаем


\begin{equation}\label{f5}


J_\beta=\frac 1{(2\pi i)^n}\sum_{\|\alpha\|\geqslant 0}\int_{\gamma(r)}


\frac{(-1)^{\|\alpha\|}\Delta\, dz}{z^{\beta+I}}\frac{Q^\alpha}{z^{


(\alpha_1+1)\beta^1+\dotsb+ (\alpha_n+1)\beta^n}}.


\end{equation}





Покажем, что в этой сумме лишь конечное число слагаемых отлично от нуля.


Для этого подсчитаем степени (по совокупности переменных) мономов, 


входящих в числитель и знаменатель подинтегрального выражения.





Степень мономов, входящих в $\Delta$ ($\deg\Delta$), не меньше


$\max(\|\beta^1\|+\dotsb+\|\beta^n\|-n,0)=\max(k_1+\dotsb+ k_n-n,0)$.


Для степени мономов, входящих в $Q^\alpha$, получаем оценку


$$


\deg Q^\alpha\geqslant \alpha_1(1+\|\beta_1\|)+\dotsb+


\alpha_n(1+\|\beta_n\|)=


\alpha_1(1+k_1)+\dotsb+\alpha_n(1+k_n).


$$


Поэтому степень числителя не меньше


$$


\max(k_1+\dotsb+ k_n-n,0)+\alpha_1(1+k_1)+\dotsb+\alpha_n(1+k_n).


$$


Степень знаменателя равна


$$


\|\beta\|+n+(\alpha_1+1)k_1+\dotsb+(\alpha_n+1)k_n.


$$


Поэтому в сумме (\ref{f5}) все слагаемые, для которых степень числителя 


больше степени знаменателя на $n$, равны нулю. Таким образом, могут 


быть отличными от нуля лишь слагаемые, для которых


$$


\max(k_1+\dotsb+ k_n-n,0)+\alpha_1(1+k_1)+\dotsb+\alpha_n(1+k_n)\leqslant


\|\beta\|+n+(\alpha_1+1)k_1+\dotsb+(\alpha_n+1)k_n.


$$


Последнее неравенство эквивалентно следующему:


$$


\|\alpha\|\leqslant \|\beta\|+\min(n,k_1+\dotsb+ k_n).\quad \square


$$





Отметим, что  в указанные в теореме~\ref{t1}


формулы, как показывает доказательство, входит лишь конечное число 


коэффициентов функций $Q_j(z)$.








\begin{cor}\label{c1}


Если все $\beta^j=(0,0,\dotsc,0)$, то


$$


J_\beta=\sum_{\|\alpha\|\leqslant\|\beta\|}(-1)^{\|\alpha\|}\frak{M}


\bigg[\frac{\Delta \, Q^\alpha}{z^\beta}\bigg]=


\sum_{\|\alpha\|\leqslant\|\beta\|}\frac{(-1)^{\|\alpha\|}}{\beta!}


\frac{\partial^{\|\beta\|}}{\partial z^\beta}(\Delta Q^\alpha)


\bigg|_{z=0}.


$$


\end{cor}





Пусть $n=1$ и $f_j(z)=f(z)=1+Q(z)=1+\sum\limits_{s=1}^\infty a_sz^s$.








\begin{cor}\label{c2} Интеграл


\begin{multline*}


J_\beta=\sum_{0\leqslant\alpha\leqslant\beta}(-1)^\alpha\frak


M\bigg(\frac{f'(z)Q^\alpha(z)}{


z^\beta}\bigg)=\sum_{0\leqslant\alpha\leqslant\beta}\frac{(-1)^\alpha}{


\beta!}\frac{\partial^\beta}{\partial z^\beta}(f'Q^\alpha)\bigg|_{z=0}= \\


%


=(-1)^{\beta+1}\begin{vmatrix} a_1 & 1 & 0 & \dotsc & 0\\


2a_2 & a_1 & 1 & \dotsc & 0\\ \hdotsfor[1.5]{5} \\


(\beta+1)a_{\beta+1} & a_\beta & a_{\beta-1} & \dotsc & a_1\end{vmatrix}.


\end{multline*}


\end{cor}





Наша дальнейшая цель --- связать рассмотренные интегралы со степенными 


суммами корней системы (\ref{f4}). Для этого нужно сузить класс  


функций $f_j$. Сначала возьмем в качестве функций $Q_j$,


$j=1,2,\dotsc, n$, многочлены вида


\begin{equation}\label{f6}


Q_j(z)=\sum_{\alpha\in M_j}a^j_\alpha z^\alpha,


\end{equation}


где $M_j$ --- конечное множество мультииндексов такое, что


$\alpha_k\leqslant \beta^j_k$, $k=1,2,\dotsc, n$,


$k\neq j$, при $\alpha\in M_j$. Но по-прежнему предполагается, что 


$\|\alpha\|>k_j$ для всех $\alpha\in M_j$.





Сделаем замену $z_j=\frac 1{w_j}$, $j=1,2,\dotsc,n$. При такой 


замене получим 


$$


f_j\bigg(\frac 1{w_1},\frac 1{w_2},\dotsc,\frac 1{w_n}\bigg)


=\frac 1{w^{\beta^j}}+


Q_j\bigg(\frac 1{w_1},\frac 1{w_2},\dotsc,\frac 1{w_n}\bigg)=


\frac 1{w^{\beta^j+s_je^j}}\bigg(w_j^{s_j}+


\widetilde Q_j(w_1,w_2,\dotsc,w_n)\bigg),


$$


где $s_j$ --- наибольшая степень многочлена $Q_j(z)$ по переменной


$z_j$,


$$


e^1=(1,0,\dotsc,0),\ e^2=(0,1,\dotsc, 0), \dotsc,\ e^n=(0,0,\dotsc, 1),


$$


а многочлен


$$


\widetilde Q_j(w_1,w_2,\dotsc,w_n)=\widetilde Q_j(w)= w^{\beta^j+s_je^j}


Q_j\bigg(\frac 1{w_1},\frac 1{w_2},\dotsc,\frac 1{w_n}\bigg)


$$


в силу условий, наложенных на многочлены $Q_j(z)$,


имеет степень по совокупности переменных строго меньшую, чем $s_j$.





По теореме Безу система нелинейных алгебраических уравнений


\begin{equation}\label{f7}


\widetilde f_j(w)=w_j^{s_j}+\widetilde Q_j(w)=0,\quad j=1,2,\dotsc, n,


\end{equation}


имеет конечное число корней, равное с учетом их кратностей


$s_1\cdot s_2\dotsb s_n$, и не имеет


корней на бесконечной гиперплоскости $\Bbb{CP}^n\setminus\Bbb C^n$.


Обозначим корни системы (\ref{f7}), не лежащие на координатных плоскостях,


с учетом их кратностей через $w_{(k)}=(w_{1(k)},w_{2(k)},\dotsc,w_{n(k)})$,


$k=1,2,\dotsc, M$, $M\leqslant s_1\cdot s_2\dotsb s_n$. Тогда точки


$z_{(k)}=\big(\frac 1{w_{1(k)}},


\frac 1{w_{2(k)}},\dotsc,\frac 1{w_{n(k)}}\big)$


и только они являются корнями системы (\ref{f4}), 


не лежащими на координатных плоскостях.








\begin{lemnonum} %%\label{l1}


Система  $(\ref{f4})$ с многочленами $f_j$ вида


$(\ref{f1})$ и многочленами $Q_j$  вида $(\ref{f6})$


имеет конечное число корней $($с учетом их кратностей\/$)$


$z_{(1)}, z_{(2)},\dotsc,z_{(M)}$, не лежащих на координатных 


плоскостях $\{z_s=0\}$, $s=1,2,\dotsc,n$.


\end{lemnonum}





Обозначим


$$


\sigma_{\beta+I}=\sigma_{(\beta_1+1,\beta_2+1,\dotsc,\beta_n+1)}=


\sum_{k=1}^M \frac 1{z_{1(k)}^{\beta_1+1}\cdot z_{2(k)}^{\beta_2+1}\dotsb


z_{n(k)}^{\beta_n+1}}.


$$


Данное выражение является степенной суммой корней, не лежащих на координатных


плоскостях, системы (\ref{f4}), но в


отрицательной степени.








\begin{thm}\label{t2}


Для системы  $(\ref{f4})$ с многочленами $f_j$ вида


$(\ref{f1})$ и многочленами $(\ref{f6})$ справедлива формула


\begin{equation}\label{f8}


J_\beta=(-1)^n\sigma_{\beta+I},


\end{equation}


где


$$


\sigma_{\beta+I}=\sum_{\|\alpha\|\leqslant\|\beta\|+\min(n,k_1+\dotsb


+k_n)}(-1)^{\|\alpha\|}


\frak{M}\bigg[\frac{\Delta\cdot Q^\alpha}{z^{\beta+(\alpha_1+1)\beta^1+


\dotsc+(\alpha_n+1)\beta^n}}\bigg].


$$


\end{thm}








\begin{pf}


Сделаем в интеграле $J_\beta$ замену переменных $z_j=\frac 1{w_j}$,


$j=1,2,\dotsc,n$. При такой замене цикл $\gamma(r)$ перейдет в цикл 


$(-1)^n\gamma\big(\frac1{r_1},\frac 1{r_2},\dotsc, \frac 1{r_n}\big)


=(-1)^n\gamma(R_1,R_2,\dotsc,R_n)$


(с учетом изменения ориентации при данной замене).





Обозначим через $\gamma^j$ мультииндекс $\beta^j+s_je^j$, $j=1,2,\dotsc,n$.


Тогда


$$


\frac{df_j\big(\frac 1{w_1},\frac 1{w_2},\dotsc,\frac 1{w_n}\big)}{


f_j\big(\frac 1{w_1},\frac 1{w_2},\dotsc,\frac 1{w_n}\big)}=


\frac{d\frac{\Tilde f_j(w)}{w^{\gamma^j}}}{


\frac{\Tilde f_j(w)}{w^{\gamma^j}}}=\frac{d \widetilde f_j(w)}{


\widetilde f_j(w) }-\sum_{k=1}^n\gamma^j_k\frac{dw_k}{w_k}.


$$


Поэтому


$$


J_\beta=\frac{(-1)^n}{(2\pi i)^n}\int_{\gamma(R)}w^{\beta+I}\bigg(


\frac{d \widetilde f_1(w)}{\widetilde f_1(w) }-\sum_{k=1}^n\gamma^1_k


\frac{dw_k}{w_k}\bigg)\land\dotsb\land\bigg(


\frac{d \widetilde f_n(w)}{\widetilde f_n(w) }-\sum_{k=1}^n\gamma^n_k


\frac{dw_k}{w_k}\bigg).


$$


Покажем, что все интегралы вида


\begin{equation}\label{f9}


\int_{\gamma(R)}w^{\beta+I} \frac{d \widetilde f_1(w)}{\widetilde f_1(w) }


\land\dotsb\land\frac{d \widetilde f_l(w)}{\widetilde f_l(w) }\land


\frac{dw_{j_1}}{w_{j_1}}\land\dotsb\land \frac{dw_{j_{n-l}}}{w_{j_{n-l}}}


\end{equation}


равны нулю, если $0\leqslant l<n$ и $R_j$ достаточно велики.





Действительно, при достаточно больших $R_j$, $j=1,2,\dotsc, n$, справедливы


неравенства


$$


|w_j|^{s_j}>|\widetilde Q_j(w)|\ \ \text{на}\ \ \gamma(R),


$$


поэтому


$$


\frac 1{\widetilde f_j(w)}=\sum_{p=0}^\infty


\frac{(-1)^p\widetilde Q_j^p(w)}{w_j^{(p+1)s_j}},


$$


следовательно, интегралы  (\ref{f9}) являются  абсолютно сходящимися


рядами из интегралов вида


$$


\int_{\gamma(R)}w^{\beta+I}\frac{w^\alpha\, dw_1\land


dw_2\land\dotsb\land dw_n}{w_1^{(p_1+1)s_1}\cdot w_2^{(p_2+1)s_2}\dotsb


w_l^{(p_l+1)s_l}\cdot w_{j_1}\dotsb w_{j_{n-l}}}.


$$


Все они равны нулю по теореме о вычетах.





Таким образом,


$$


J_\beta=\frac{(-1)^n}{(2\pi i)^n}\int_{\gamma(R)}w^{\beta+I}


\frac{d \widetilde f_1(w)}{\widetilde f_1(w) }


\land\dotsb\land\frac{d \widetilde f_n(w)}{\widetilde f_n(w) }.


$$





По теореме Южакова (напр., \cite{1}, гл.\,1) о логарифмическом вычете


последний интеграл равен сумме значений голоморфной функции $w^{\beta+I}$


во всех корнях системы (\ref{f7}). Но значение функции $w^{\beta+I}$ 


в корне системы (\ref{f7}), лежащем на координатной плоскости, равно нулю.


Поэтому выполнено \eqref{f8}.


\end{pf}








\begin{exam}


Рассмотрим систему


$$


f_j(z)=1+a_jz_j=0, \quad j=1,2,\dotsc, n.


$$


Тогда $\Delta=a_1\cdot a_2\dotsb a_n$. Единственный корень имеет координаты


$z_j=-\frac 1{a_j}$, $j=1,2,\dotsc, n$, и


$$


\sigma_{\beta+I}=(-1)^{\|\beta\|+n} a_1^{\beta_1+1}\dotsb


a_n^{\beta_n+1}=(-1)^{\|\beta\|+n}a^{\beta+I}.


$$


Многочлены $Q^\alpha(z)=a^\alpha z^\alpha$,


а


$\beta^j=(0,0,\dotsc,0)$, поэтому


$$


J_\beta=\sum_{\|\alpha\|\leqslant\|\beta\|}(-1)^{\|\alpha\|}\frak{M}


\bigg[\frac{\Delta \, Q^\alpha}{z^\beta}\bigg]=


(-1)^{\|\beta\|}a_1\dotsb a_n\cdot a^\beta=(-1)^{\|\beta\|} a^{\beta +I}


$$


в полном соответствии с теоремой \ref{t2}.


\end{exam}





Рассмотрим более общую ситуацию. Пусть


\begin{equation}\label{f10}


f_j(z)=\frac{f^{(1)}_j(z)}{f^{(2)}_j(z)},\quad j=1,2,\dotsc, n,


\end{equation}


где $f^{(1)}_j(z)$, $f^{(2)}_j(z)$ --- целые функции в $\Bbb C^n$, 


разлагающиеся в бесконечные произведения (равномерно и абсолютно 


сходящиеся в $\Bbb C^n$)


$$


f^{(1)}_j(z)=\prod_{s=1}^\infty f^{(1)}_{js}(z),\quad f^{(2)}_{j}(z)=


\prod_{s=1}^\infty f^{(2)}_{js}(z),


$$


причем каждый из сомножителей имеет форму $z^{\beta^j}+Q_j(z)$,


а $Q_j(z)$ --- функции вида (\ref{f6}).





Для каждого набора индексов $j_1,\dotsc,j_n\in\Bbb N$


и каждого набора чисел $i_1,\dotsc, i_{n}$, равных


1~или~2, системы нелинейных алгебраических уравнений


\begin{equation}\label{f11}


f^{(i_1)}_{1j_1}(z)=0,\ \ f^{(i_2)}_{2j_2}(z)=0,\ \dotsc,


\ f^{(i_n)}_{nj_n}(z)=0


\end{equation}


имеют согласно лемме конечное число корней, не лежащих на


координатных плоскостях.





Корни всех таких систем, не лежащие на координатных плоскостях, составляют 


не более чем счетное множество. Перенумеруем


их с учетом кратностей: $z_{(1)},\  z_{(2)},\dotsc, z_{(l)},\dotsc$


Будем предполагать, что ряд


\begin{equation}\label{f12}


\sum_{l=1}^\infty \frac 1{|z_{1(l)}|\cdot |z_{2(l)}|\dotsb|z_{n(l)}|}


\end{equation}


сходится. Обозначим 


$$


\sigma_{\beta+I}=\sum_{l=1}^\infty \frac{\varepsilon_l}{z^{\beta_1+1}_{1(l)}


\cdot z^{\beta_2+1}_{2(l)}\dotsb z^{\beta_n+1}_{n(l)}}.


$$


Здесь $\beta_1,\dotsc,\beta_n$, как и прежде, --- неотрицательные целые числа,


а знак $\varepsilon_l$ равен $+1$, если в систему (\ref{f11}), корнем


которой является $z_{(l)}$, входит четное число функций $f^{(2)}_{js}$; и


равен $-1$, если входит нечетное число функций $f^{(2)}_{js}$.





Ряды, определяющие суммы $\sigma_{\beta+I}$, сходятся в силу условия,


наложенного на ряд (\ref{f12}).





Для системы (\ref{f4}), составленной из функций (\ref{f10}), точки


$z_{(l)}$ являются корнями или особыми точками (полюсами). 


Модули $|z_{(l)}|$ не могут стремиться к нулю. В противном случае


ряд (\ref{f12}) не мог бы сходиться. Следовательно, все функции


$f_j$ голоморфны и не равны нулю


вблизи точки~0 и вне координатных плоскостей, поэтому для них определены


интегралы $J_\beta$.








\begin{thm}\label{t3}


Для системы $(\ref{f4})$ с функциями $(\ref{f10})$


справедливы равенства \eqref{f8}.


\end{thm}








\begin{pf}


Так как


$$


\frac{df_j(z)}{f_j(z)}=\frac{d\,f^{(1)}_j(z)} {f^{(1)}_j(z)}


-\frac{d\,f^{(2)}_j(z)} {f^{(2)}_j(z)},


$$


то


\begin{multline*}


\frac{df_1(z)}{f_1(z)}\land \frac{df_2(z)}{f_2(z)}


\land\dotsb\land \frac{df_n(z)}{f_n(z)}= \\


%


=\bigg(\frac{d\,f^{(1)}_1(z)} {f^{(1)}_1(z)}-\frac{d\,f^{(2)}_1(z)}


{f^{(2)}_1(z)}\bigg)


\land \bigg(\frac{d\,f^{(1)}_2(z)} {f^{(1)}_2(z)}-\frac{d\,f^{(2)}_2(z)}


{f^{(2)}_2(z)}\bigg)\land


\dotsb\land \bigg(\frac{d\,f^{(1)}_n(z)} {f^{(1)}_n(z)}-


\frac{d\,f^{(2)}_n(z)} {f^{(2)}_n(z)}\bigg)= \\


%


=\sum (-1)^s \frac{d\,f^{(i_1)}_1(z)} {f^{(i_1)}_1(z)}\land


\frac{d\,f^{(i_2)}_2(z)} {f^{(i_2)}_2(z)}\land\dotsb\land


\frac{d\,f^{(i_n)}_n(z)}{f^{(i_n)}_n(z)},


\end{multline*}


где $s$ --- число сомножителей, для которых $i_l=2$, а сумма берется


по всевозможным наборам чисел $i_1, i_2,\dotsc, i_n$, равных 1 или 2.





Поэтому теорему достаточно доказать для целых функций $f_j(z)$. 


В этом случае


$$


\frac{d f_j(z)}{f_j(z)}=\frac{d\prod\limits_{s=1}^\infty f_{js}(z)}{


\prod\limits_{s=1}^\infty f_{js}(z)}=\sum_{s=1}^\infty


\frac{d f_{js}(z)}{f_{js}(z)}.


$$


Последний ряд сходится абсолютно и равномерно на $\gamma(r)$ для достаточно


малых $r_j$. Таким образом, интеграл $J_\beta$ равен сумме интегралов вида


$$


\frac 1{(2\pi i)^n}\int_{\gamma(r)}\frac 1{z^{\beta+I}}


\frac{d f_{1s_1}(z)}{f_{1s_1}(z)}\land\frac{d\, f_{2s_2}(z)}{f_{2s_2}(z)}


\land\dotsb\land \frac{d\, f_{ns_n}(z)}{f_{ns_n}(z)}.


$$


Для каждого из этих интегралов нужная формула доказана (теорема \ref{t2}).


\end{pf}





Отметим, что если $f_j(z)$ являются целыми функциями, разлагающимися в


бесконечные произведения, то они сами имеют вид (\ref{f1}) с функциями $Q_j(z)$


вида (\ref{f2}). Поэтому интегралы $J_\beta$ вычисляются по теореме~\ref{t1}.








\section{Приложения}








\begin{exam}


Рассмотрим систему уравнений от двух комплексных переменных


\begin{equation}\label{f13}


\begin{split} f_1(z_1,z_2)&=1+a_1z_1=0,\\


f_2(z_1,z_2)&=1+b_1z_1+b_2z_2=0. \end{split}


\end{equation}


Здесь функции не удовлетворяют условиям теоремы \ref{t2}, но удовлетворяют 


условиям теоремы \ref{t1}. Поэтому


\begin{gather*}


J_\beta=\frac 1{(2\pi i)^2}\int_{\gamma(r)}\frac


1{z_1^{\beta_1+1}z_2^{\beta_2+1}}


\frac{df_1\land df_2}{f_1f_2}=


\frac 1{(2\pi i)^2}\int_{\gamma(r)}\frac 1{z_1^{\beta_1+1}z_2^{\beta_2+1}}


\frac{a_1b_2 dz_1\land dz_2}{(1+a_1z_1)(1+b_1z_1+b_2z_2)}= \\


%


=\frac{a_1b_2}{\beta_1!\cdot\beta_2!}\frac{\partial^{\beta_1+\beta_2}}{


\partial z_1^{\beta_1}\partial z_2^{\beta_2}}\bigg[\frac 1


{(1+a_1z_1)(1+b_1z_1+b_2z_2)}\bigg]_{z_1=z_2=0}= \\


%


=\frac{a_1b_2}{\beta_1!} \frac{\partial^{\beta_1}}{


\partial z_1^{\beta_1}}\bigg[\frac 1{(1+a_1z_1)}\frac{(-1)^{\beta_2}


b_2^{\beta_2}}{(1+b_1z_1+b_2z_2)^{\beta_2+1}}\bigg]_{z_1=z_2=0}.


\end{gather*}


Применяя формулу Лейбница дифференцирования произведения 


двух функций, получим


\begin{gather*}


J_\beta=\frac{(-1)^{\beta_2}a_1b_2^{\beta_2+1}}{\beta_1!}


\sum_{s=0}^{\beta_1}\frac{\beta_1!}{s! (\beta_1-s)!}


\bigg[\frac{\partial^s}{\partial z_1^s}\bigg(


\frac 1{1+a_1z_1}\bigg)\frac{\partial^{\beta_1-s}}


{\partial z_1^{\beta_1-s}}


\bigg(\frac 1{(1+b_1z_1)^{\beta_2+1}}\bigg)\bigg]_{z_1=0}= \\


%


=(-1)^{\beta_2}a_1b_2^{\beta_2+1}\sum_{s=0}^{\beta_1}


\frac{1}{s! (\beta_1-s)!}


\bigg[\frac{(-1)^s s! a_1^s}{(1+a_1z_1)^{s+1}}


\frac{(-1)^{\beta_1-s} b_1^{\beta_1-s} (\beta_2+1)\dotsb(\beta_1+\beta_2-s)}


{(1+b_1z_1)^{\beta_1+\beta_2+1-s}}\bigg]_{z_1=0}= \\


%


=(-1)^{\beta_1+\beta_2}a_1b_2^{\beta_2+1}


\sum_{s=0}^{\beta_1}\frac{a_1^sb_1^{\beta_1-s}


(\beta_1+\beta_2-s)!}{(\beta_1-s)!\beta_2!}= \\


%


=\frac{(-1)^{\beta_1+\beta_2}a_1b_2^{\beta_2+1}}{\beta_2!}


\sum_{s=0}^{\beta_1}\frac{(\beta_1+\beta_2-s)!}{(\beta_1-s)!}


a_1^sb_1^{\beta_1-s}.


\end{gather*}


Например, в случае $\beta_2=0$ имеем


$$


J_{(\beta_1,0)}=(-1)^{\beta_1}a_1b_2\frac{a_1^{\beta_1+1}-b_1^{\beta_1+1}}


{a_1-b_1}.


$$





Корень системы (\ref{f13}) равен $z_1=-\frac 1{a_1}$,


$z_2=\frac{b_1-a_1}{a_1b_2}$. Поэтому степенная сумма


$$


\sigma_{\beta+I}=\frac{(-1)^{\beta_1+\beta_2}a_1^{\beta_1+\beta_2+2}


b_2^{\beta_2+1}}{(a_1-b_1)^{\beta_2+1}},


$$


в частности,


$$


\sigma_{(\beta_1+1,1)}=\frac{(-1)^{\beta_1}a_1^{\beta_1+2}b_2}{a_1-b_1},


$$


т.\,е.


\begin{equation}\label{13'}


I_{(\beta_1,0)}=


\sigma_{(\beta_1+1,1)}-\frac{(-1)^{\beta_1}a_1b_2b_1^{\beta_1+1}}{a_1-b_1}.


\end{equation}


Значит, теорема \ref{t2} для системы (\ref{f13}) не верна (при $b_2\neq 0$).


Тем не менее проведенные вычисления и теорема \ref{t2} позволяют находить


суммы некоторых двойных рядов.


\end{exam}





Напомним известные разложения синуса в бесконечное произведение и степенной


ряд:


$$


\frac{\sin\sqrt z}{\sqrt z}=\prod_{k=1}^\infty


\bigg(1-\frac z{k^2\pi^2}\bigg)=


\sum_{k=0}^\infty \frac{(-1)^kz^k}{(2k+1)!},


$$


которые равномерно и абсолютно сходятся на комплексной плоскости.





Рассмотрим систему уравнений


\begin{equation}\label{f14}


\begin{split}


f_1(z_1,z_2)&=\dfrac{\sin\sqrt{z_1}}{\sqrt{z_1}}=


\prod\limits_{k=1}^\infty\bigg(1-\dfrac{z_1}{k^2\pi^2}\bigg)=0,\\


f_2(z_1,z_2)&=\dfrac{\sin\sqrt{z_2-z_1}}{\sqrt{z_2-z_1}}=


\prod\limits_{m=1}^\infty\bigg(


1-\dfrac{z_2-z_1}{m^2\pi^2}\bigg)=0. \end{split}


\end{equation}





Каждая из функций этой системы разлагается в бесконечное произведение 


функций из системы (\ref{f13}). Корнями системы (\ref{f14}) являются точки


$(\pi^2k^2,\pi^2(k^2+m^2))$, $k,m\in\Bbb N$. Поэтому степенная сумма


$\sigma_{(\beta_1,\beta_2)}$ равна сумме ряда:


$$


\sigma_{(\beta_1,\beta_2)}=\frac 1{\pi^{2(\beta_1+\beta_2)}}


\sum_{k,m=1}^\infty


\frac 1{k^{2\beta_1}(k^2+m^2)^{2\beta_2}}.


$$


Для системы (\ref{f14}) верна теорема \ref{t1}, поскольку


$$


f_1=\sum_{k=0}^\infty\frac {(-1)^k z_1^k}{(2k+1)!},


$$


а


$$


f_2=\sum_{k=0}^\infty\frac {(-1)^k (z_2-z_1)^k}{(2k+1)!}=


\sum_{m,n=0}^\infty\frac {(-1)^m (m+n)! z_1^n z_2^m}{m!n!(2m+2n+1)!}.


$$





Рассмотрим интеграл $J_\beta$ для системы (\ref{f14}).


Используя равенство (\ref{13'}) и вид корней системы (\ref{f14}), 


получим


$$


J_{(\beta_1,0)}=\sigma_{(\beta_1+1,1)}+(-1)^{\beta_1}


\sum_{k,m=1}^\infty\frac


1{ \pi^{2(\beta_1+1)}(k^2+m^2)m^{2(\beta_1+1)}}.


$$


Таким образом, для нечетных $\beta_1$ интеграл $I_{(\beta_1,0)}=0$, а для


четных $\beta_1$ интеграл $I_{(\beta_1,0)}=2\sigma_{(\beta_1+1,1)}$.


Полагая $\beta_1=2s$, получим следующую формулу для нахождения суммы ряда.








\begin{cor}\label{c3} 


Справедливы формулы


\begin{equation*} %%%\label{f15}


\sigma_{(2s+1,1)}=\frac 1{\pi^{2(2s+1)}}


\sum_{k,m=1}^\infty \frac 1{k^{2(2s+1)}


(m^2+k^2)}=\frac 12 J_{(2s,0)}=\frac 12\sum_{\|\alpha\|\leqslant 2s}


\frak M\bigg[\frac{\Delta Q^\alpha}{z_1^{2s}}\bigg],


\end{equation*}


где $\Delta$, $Q$ и $\frak M$ определены в теореме~\ref{t1}.


\end{cor}





Вычислим, например, $I_{(2,0)}$. Якобиан


$$


\Delta=\frac{\partial f_1}{\partial z_1}


\frac{\partial f_2}{\partial z_2}=\frac 1{(3!)^2}-\frac{2z_2}{3!5!}


+\frac{4z_1z_2}{(5!)^2}+ z_1^2


\bigg(\frac 1{7!}-\frac 4{(5!)^2}\bigg)+\dotsb,


$$


тогда


\begin{gather*}


\frak M\bigg[\frac\Delta{z_1^2}\bigg]=\frac 1{7!}-\frac 4{(5!)^2},\quad


\frak M\bigg[\frac{\Delta Q_1}{z_1^2}\bigg]=\frac 1{(3!)^25!},\quad


\frak M\bigg[\frac{\Delta Q_2}{z_1^2}\bigg]=\frac 1{(3!)^25!}, \\


%


\frak M\bigg[\frac{\Delta Q_1^2}{z_1^2}\bigg]= \frac 1{(3!)^4},\quad


\frak M\bigg[\frac{\Delta Q_1\cdot Q_2}{z_1^2}\bigg]=-\frac 1{(3!)^4},\quad


\frak M\bigg[\frac{\Delta Q_2^2}{z_1^2}\bigg]=\frac 1{(3!)^4}.


\end{gather*}


Поэтому  $J_{(2,0)}=\frac{13}{56700}$,


а $\sigma_{(3,1)}=\frac{13}{113400}$. Следовательно,


$$


\sum_{k,m=1}^\infty\frac 1{k^6(k^2+m^2)}=\frac{13\pi^8}{113400},


$$


что отличается от формулы 10 (\cite{6}, с.\,750)


(где она приведена с ошибкой),


но полностью совпадает с примером 1 из \cite{7}.





Рассмотрим более сложную систему


\begin{equation}\label{f16}


\begin{split}


f_1(z_1,z_2)&=\dfrac{\sin\sqrt{z_1-a^2}}{\sqrt{z_1-a^2}}=


\prod\limits_{k=1}^\infty


\bigg(1-\dfrac{z_1-a^2}{\pi^2k^2}\bigg)=0,\\


f_1(z_1,z_2)&=\dfrac{\sin\sqrt{z_2-z_1-a^2}}{\sqrt{z_2-z_1-a^2}}=


\prod\limits_{m=1}^\infty


\bigg(1-\dfrac{z_2-z_1-a^2}{\pi^2m^2}\bigg)=0.\end{split}


\end{equation}


Корни данной системы равны $(\pi^2k^2+a^2,\pi^2(m^2+k^2)+2a^2)$,


$k,m\in\Bbb N$. Функции системы (\ref{f16}) являются бесконечными 


произведениями функций из системы (\ref{f13}).  Данная система 


удовлетворяет теореме \ref{t1}. Нетрудно подсчитать (используем, 


как для предыдущей системы~(\ref{f14}), формулу~(\ref{13'})), что


интеграл $J_{(0,0)}=2\sigma_{(1,1)}.$





Для функций


$$


f_1=\sum_{k=0}^\infty\frac{(-1)^k(z_1-a^2)^k}{(2k+1)!},\quad


f_2=\sum_{k=0}^\infty\frac{(-1)^k(z_2-z_1-a^2)^k}{(2k+1)!}


$$


имеем


$$


f_1(0,0)=f_2(0,0)=\sum_{k=0}^\infty\frac {a^{2k}}{(2k+1)!}=\frac{\sh a}a.


$$


Следовательно, чтобы применить формулу из теоремы~\ref{t1}, нужно функции


$f_1$ и $f_2$ разделить на эту константу (отнормировать).





Тогда интеграл


$$


J_{(0,0)}=\frak M [\Delta]=\bigg(\frac 1{2a^2}-\frac 1{2a}\cth a\bigg)^2,


$$


отсюда получим


$$


\sigma_{(1,1)}=\sum_{k,m=1}^\infty\frac 1{\pi^2k^2+a^2}


\frac1{\pi^2(k^2+m^2)+2a^2}=\frac{(1-a\cth a)^2}{8a^4}.


$$








\begin{thebibliography}{9}





\bibitem{1}


Айзенберг Л.А., Южаков А.П. {\it Интегральные представления и


вычеты в многомерном комплексном анализе}. -- Новосибирск: Наука,


1979. -- 366~c.


\bibitem{2}


Цих А.К. {\it Многомерные вычеты и их применения}. --


Новосибирск: Наука, 1989. -- 240~c.


\bibitem{3}


Быков В.И., Кытманов  А.М., Лазман М.З. {\it Методы исключения  в 


компьютерной алгебре многочленов.} -- Новосибирск: Наука, 1991.


-- 234~c.


\bibitem{4}


Быков В.И., Кытманов  А.М., Осетрова Т.А. {\it Компьютерная


алгебра многочленов. Модифицированный метод исключения} //


Докл. РАН. -- 1996. -- Т.\,350. -- \No\,4. -- С.\,443--445.


\bibitem{5}


Быков В.И., Кытманов  А.М., Осетрова Т.А., Потапова З.Е.


{\it Применение систем компьютерной алгебры в модифицированном методе 


исключения неизвестных} //


Докл. РАН. -- 2000. -- Т.\,370. -- \No\,4. -- С.\,439--442.


\bibitem{6}


Прудников А.П., Брычков Ю.А., Маричев О.И.


{\it Интегралы и ряды. Элементарные функции.} -- М.: Наука, 1981. -- 800~с.


\bibitem{7}


Ермилов И.В. {\it Вычисление сумм и улучшение сходимости числовых рядов}


// Исследов. по комплексному


анализу. -- Изд-во Красноярск. ун-та, 1989. -- С.\,52--62.





\end{thebibliography}





\vskip 2cm





\post{\it Красноярский государственный университет}{\it Поступила}


\post{\it Московский авиационный институт}{07.04.2003}





\label{end}


\end{document}


